\chapter{Backend: reportes, votacion y roles}
\label{ch:backend_reportes}

\section{Objetivo}

Extender el backend con:

\begin{itemize}[leftmargin=*]
  \item Sistema de reportes para publicaciones y canales.
  \item Votacion colegiada de administradores (3 votos para resolver).
  \item Gestion de roles con jerarquia (STUDENT, PROFESSOR, ADMIN, SUPERADMIN).
  \item Baneo y desbaneo de usuarios.
  \item Log de auditoria de acciones administrativas.
  \item Moderacion directa por dueno del canal.
\end{itemize}

\section{Contexto}

Hasta la Sesion~\ref{ch:sesion5}, el backend tenia el dominio basico: usuarios, canales, suscripciones, publicaciones y notificaciones. No existia ningun mecanismo de moderacion ni distincion real entre roles mas alla del valor del enum.

Se decidio introducir tres reglas clave:

\begin{enumerate}[leftmargin=*]
  \item Solo el SUPERADMIN puede promover/degradar libremente (excepto quitar SUPERADMIN).
  \item Los ADMIN solo pueden promover estudiantes a PROFESSOR o ADMIN.
  \item Los duenos de un canal pueden moderar sus publicaciones sin votacion.
\end{enumerate}

\section{Desarrollo}

\subsection{Migracion V2}

Se creo \texttt{V2\_\_reports\_and\_roles.sql} con tres tablas:

\begin{itemize}[leftmargin=*]
  \item \texttt{reports}: un reporte por (target, reporter).
  \item \texttt{report\_votes}: un voto por (report, admin).
  \item \texttt{admin\_action\_log}: bitacora de acciones administrativas.
\end{itemize}

\begin{lstlisting}[language=SQL]
CREATE TABLE reports (
    id           UUID PRIMARY KEY DEFAULT gen_random_uuid(),
    target_type  VARCHAR(20)  NOT NULL CHECK (target_type IN ('POST', 'CHANNEL')),
    target_id    UUID         NOT NULL,
    reporter_id  UUID         NOT NULL REFERENCES users(id) ON DELETE CASCADE,
    reason       VARCHAR(50)  NOT NULL,
    details      VARCHAR(500),
    status       VARCHAR(30)  NOT NULL DEFAULT 'PENDING',
    created_at   TIMESTAMPTZ  NOT NULL DEFAULT NOW(),
    resolved_at  TIMESTAMPTZ,
    resolved_by  UUID REFERENCES users(id) ON DELETE SET NULL,
    UNIQUE (target_type, target_id, reporter_id)
);

CREATE TABLE report_votes (
    id          UUID PRIMARY KEY DEFAULT gen_random_uuid(),
    report_id   UUID NOT NULL REFERENCES reports(id) ON DELETE CASCADE,
    admin_id    UUID NOT NULL REFERENCES users(id) ON DELETE CASCADE,
    decision    VARCHAR(20) NOT NULL CHECK (decision IN ('KEEP', 'DELETE')),
    comment     VARCHAR(500),
    voted_at    TIMESTAMPTZ NOT NULL DEFAULT NOW(),
    UNIQUE (report_id, admin_id)
);

CREATE TABLE admin_action_log (
    id           UUID PRIMARY KEY DEFAULT gen_random_uuid(),
    admin_id     UUID NOT NULL REFERENCES users(id) ON DELETE CASCADE,
    action       VARCHAR(50) NOT NULL,
    target_type  VARCHAR(20),
    target_id    UUID,
    details      VARCHAR(500),
    created_at   TIMESTAMPTZ NOT NULL DEFAULT NOW()
);
\end{lstlisting}

\subsection{Enumeraciones}

\begin{lstlisting}[language=Java]
public enum Role { STUDENT, PROFESSOR, ADMIN, SUPERADMIN }
public enum ReportTargetType { POST, CHANNEL }
public enum ReportReason { SPAM, INAPPROPRIATE, MISINFORMATION, HARASSMENT, OTHER }
public enum ReportStatus { PENDING, RESOLVED_KEEP, RESOLVED_DELETE }
public enum VoteDecision { KEEP, DELETE }
\end{lstlisting}

\subsection{Entidades JPA}

Se crearon \texttt{Report}, \texttt{ReportVote} y \texttt{AdminActionLog} siguiendo el mismo patron que las entidades previas: UUID como clave primaria, \texttt{@CreatedDate} para auditoria.

\subsection{Repositorios}

\begin{lstlisting}[language=Java]
public interface ReportRepository extends JpaRepository<Report, UUID> {
    Page<Report> findByStatusOrderByCreatedAtDesc(ReportStatus status, Pageable p);
    Page<Report> findAllByOrderByCreatedAtDesc(Pageable p);
    List<Report> findByTargetTypeAndTargetIdAndStatus(
        ReportTargetType type, UUID targetId, ReportStatus status);
    long countByStatus(ReportStatus status);
    boolean existsByTargetTypeAndTargetIdAndReporterId(
        ReportTargetType type, UUID targetId, UUID reporterId);
}

public interface ReportVoteRepository extends JpaRepository<ReportVote, UUID> {
    List<ReportVote> findByReportId(UUID reportId);
    Optional<ReportVote> findByReportIdAndAdminId(UUID reportId, UUID adminId);
    long countByReportIdAndDecision(UUID reportId, VoteDecision decision);
    long countByReportId(UUID reportId);
}

public interface AdminActionLogRepository extends JpaRepository<AdminActionLog, UUID> {
    Page<AdminActionLog> findAllByOrderByCreatedAtDesc(Pageable p);
    Page<AdminActionLog> findByAdminIdOrderByCreatedAtDesc(UUID adminId, Pageable p);
}
\end{lstlisting}

\subsection{\texttt{ReportService}}

Es el nucleo del sistema de moderacion colegiada. Reglas:

\begin{enumerate}[leftmargin=*]
  \item Cualquier usuario autenticado puede reportar.
  \item Un mismo usuario no puede reportar dos veces el mismo contenido.
  \item Un admin no puede votar en un reporte que el mismo creo.
  \item Cada admin vota una sola vez por reporte.
  \item Al alcanzar 3 votos, se resuelve por mayoria simple.
  \item Si gana DELETE, el contenido se elimina automaticamente.
\end{enumerate}

\begin{lstlisting}[language=Java]
public static final int VOTES_REQUIRED = 3;

private void resolveIfComplete(Report report) {
    long totalVotes = voteRepository.countByReportId(report.getId());
    if (totalVotes < VOTES_REQUIRED) return;

    long deleteVotes = voteRepository.countByReportIdAndDecision(
        report.getId(), VoteDecision.DELETE);
    long keepVotes = voteRepository.countByReportIdAndDecision(
        report.getId(), VoteDecision.KEEP);

    ReportStatus newStatus;
    if (deleteVotes > keepVotes) {
        newStatus = ReportStatus.RESOLVED_DELETE;
        applyDeletion(report);
    } else {
        newStatus = ReportStatus.RESOLVED_KEEP;
    }

    report.setStatus(newStatus);
    report.setResolvedAt(Instant.now());
    reportRepository.save(report);
}
\end{lstlisting}

\subsection{\texttt{AdminUserService}}

Encapsula la logica de cambio de rol y ban. Reglas aplicadas en orden:

\begin{enumerate}[leftmargin=*]
  \item Nadie puede cambiarse su propio rol.
  \item Nunca se puede quitar el rol SUPERADMIN.
  \item Solo SUPERADMIN puede nombrar a otro SUPERADMIN.
  \item ADMIN solo puede promover STUDENT a PROFESSOR o ADMIN.
  \item Solo SUPERADMIN puede degradar.
\end{enumerate}

\begin{lstlisting}[language=Java]
public UserSummaryResponse changeRole(UUID targetUserId, Role newRole) {
    var me = currentUser.require();
    UUID myId = me.id();
    Role myRole = Role.valueOf(me.role());

    if (myId.equals(targetUserId))
        throw new ForbiddenException("No puedes cambiar tu propio rol");

    User target = userRepository.findById(targetUserId)
        .orElseThrow(() -> new NotFoundException("Usuario no encontrado"));

    Role currentRole = target.getRole();

    if (currentRole == Role.SUPERADMIN && newRole != Role.SUPERADMIN)
        throw new ForbiddenException("No se puede degradar a un SUPERADMIN");

    if (newRole == Role.SUPERADMIN && myRole != Role.SUPERADMIN)
        throw new ForbiddenException("Solo un SUPERADMIN puede nombrar a otro SUPERADMIN");

    if (myRole == Role.ADMIN) {
        if (currentRole != Role.STUDENT)
            throw new ForbiddenException("Como ADMIN solo puedes promover estudiantes");
        if (newRole != Role.PROFESSOR && newRole != Role.ADMIN)
            throw new ForbiddenException("Como ADMIN solo puedes promover a PROFESSOR o ADMIN");
    }

    if (myRole != Role.SUPERADMIN && level(newRole) < level(currentRole))
        throw new ForbiddenException("Solo un SUPERADMIN puede degradar usuarios");

    target.setRole(newRole);
    userRepository.save(target);
    // ...
    return toResponse(target);
}
\end{lstlisting}

\subsection{Moderacion directa por dueno del canal}

En \texttt{PostService.delete()} se anadio la regla que permite al dueno del canal (incluyendo profesores) borrar publicaciones ajenas sin necesidad de votacion.

\begin{lstlisting}[language=Java]
@Transactional
public void delete(UUID postId) {
    Post p = postRepository.findById(postId)
        .orElseThrow(() -> new NotFoundException("Publicacion no encontrada"));

    var me = currentUser.require();
    UUID myId = me.id();
    Role myRole = Role.valueOf(me.role());

    UUID authorId = p.getAuthor().getId();
    UUID channelOwnerId = p.getChannel().getOwner().getId();

    boolean isAuthor = authorId.equals(myId);
    boolean isChannelOwner = channelOwnerId.equals(myId);
    boolean isAdmin = myRole == Role.ADMIN || myRole == Role.SUPERADMIN;

    if (!isAuthor && !isChannelOwner && !isAdmin) {
        throw new ForbiddenException(
            "Solo el autor, el dueno del canal o un moderador pueden eliminar");
    }

    if (!isAuthor) {
        User admin = userRepository.findById(myId).orElseThrow();
        String reason = isChannelOwner
            ? "Moderacion del dueno del canal"
            : "Moderacion administrativa";
        adminLog.log(admin, "DELETE_POST", ReportTargetType.POST,
            postId, reason + ": " + p.getChannel().getName());
    }

    postRepository.delete(p);
}
\end{lstlisting}

\subsection{Endpoints}

\begin{lstlisting}[language=bash]
POST   /api/reports                          (cualquier usuario)
GET    /api/admin/reports?status=...         (ADMIN, SUPERADMIN)
GET    /api/admin/reports/{id}               (ADMIN, SUPERADMIN)
POST   /api/admin/reports/{id}/vote          (ADMIN, SUPERADMIN)
GET    /api/admin/users?q=...                (ADMIN, SUPERADMIN)
GET    /api/admin/users/{id}                 (ADMIN, SUPERADMIN)
PATCH  /api/admin/users/{id}/role            (ADMIN, SUPERADMIN con reglas)
PATCH  /api/admin/users/{id}/status          (solo SUPERADMIN)
GET    /api/admin/audit                      (ADMIN propio / SUPERADMIN todo)
DELETE /api/posts/{id}                       (autor, dueno, moderador)
\end{lstlisting}

\subsection{Habilitacion de \texttt{@PreAuthorize}}

Se agrego \texttt{@EnableMethodSecurity} a \texttt{SecurityConfig} para que las anotaciones de metodo sean respetadas.

\subsection{Problemas encontrados y solucion}

\subsubsection{\texttt{LazyInitializationException} en el endpoint de auditoria}

El endpoint \texttt{/api/admin/audit} devolvia 403 en lugar de 200. El log revelo una excepcion oculta:

\begin{lstlisting}[language=bash]
org.hibernate.LazyInitializationException: Could not initialize proxy
[mx.uam.notiuam.domain.User#...] - no session
    at ...AdminAuditController.lambda$list$0(AdminAuditController.java:41)
\end{lstlisting}

El mapeo de entidades a DTO se hacia en el controller, fuera de la transaccion, y al acceder a \texttt{log.getAdmin().getDisplayName()} la sesion JPA ya estaba cerrada.

\textbf{Solucion:} mover el mapeo al servicio (\texttt{AdminLogService}), dentro del \texttt{@Transactional(readOnly = true)}.

\section{Verificacion}

\subsection{Reglas de rol}

\begin{itemize}[leftmargin=*]
  \item SUPERADMIN promueve y degrada sin restricciones (salvo quitar SUPERADMIN).
  \item ADMIN solo promueve estudiantes.
  \item ADMIN no puede banear.
  \item Un usuario baneado no puede iniciar sesion.
\end{itemize}

\subsection{Flujo de reporte y votacion}

\begin{enumerate}[leftmargin=*]
  \item Betzabet reporta un post de Aaron.
  \item Guillermo vota DELETE.
  \item Aaron vota DELETE.
  \item Un cuarto admin vota KEEP.
  \item Con 3 votos (2--1) el sistema resuelve con \texttt{RESOLVED\_DELETE}.
  \item El post se elimina automaticamente.
\end{enumerate}

\subsection{Moderacion directa}

\begin{itemize}[leftmargin=*]
  \item El dueno del canal puede borrar publicaciones ajenas.
  \item El borrado se registra como \texttt{Moderacion del dueno del canal}.
  \item Un ADMIN puede borrar en cualquier canal.
  \item Un STUDENT solo puede borrar sus propios posts.
\end{itemize}

\section{Resultado}

\begin{exito}[Estado al cierre]
\begin{itemize}[leftmargin=*]
  \item Sistema de moderacion completo con votacion colegiada.
  \item Jerarquia de roles con reglas bien definidas.
  \item Auditoria de todas las acciones administrativas.
  \item Base solida para la interfaz de administracion en el cliente Flutter.
\end{itemize}
\end{exito}